\section{Formulación variacional de Einstein--Cartan: restricciones
hamiltonianas y evolución Einstein--Schrödinger}
\label{sec:gravedad-einstein-schrodinger-hamiltoniana-rev11}
\index{Einstein--Schrödinger!ecuación HMT}
\index{gravedad!formulación hamiltoniana}

La acción de primer orden y la ecuación métrica poseen una formulación
hamiltoniana restringida. Su realización cuántica se tipa mediante un espacio
de secciones \(\mathcal H_{\rm sec}\), un dominio denso invariante
\(\mathscr D\subset\mathcal H_{\rm sec}\), un mapa de cuantización
\(\mathcal Q\) y un ordenamiento de las restricciones que preserve
\(\mathscr D\). Sean \(g^{(A)}\) y \(g^{(V)}\) las métricas areal y
volumétrica, y sea \(\theta_T\) la coordenada torsional. Fijados esos datos,
el funcional de estado \(\Psi\in\mathscr D\) satisface
\begin{equation}
 i\hbar_{\rm HMT}\,\partial_t\Psi
 =\widehat H_{\rm HMT}
 [g^{(A)},g^{(V)},\theta_T;\gamma_{\rm HMT},R_5(q_{AC})]\Psi,
 \label{eq:einstein-schrodinger-hmt-rev11}
\end{equation}
con la restricción geométrica
\begin{equation}
 \left(
 G_{\mu\nu}+\Lambda_{\rm eff}g_{\mu\nu}
 +U_{\mu\nu}(T)
 -\frac{8\pi G_{\rm HMT}}{c_{\rm int}^4}
 \widehat T_{\mu\nu}
 \right)\Psi=0.
 \label{eq:restriccion-einstein-schrodinger-hmt-rev11}
\end{equation}
La primera ecuación gobierna la evolución de la sección; la segunda impone
la compatibilidad de la misma sección con curvatura, torsión y materia. Son,
respectivamente, la ecuación dinámica y la familia de restricciones de una
única realización sobre \(\mathscr D\). Su consistencia operatoria significa
que la evolución preserva \(\mathscr D\) y que los conmutadores de las
restricciones cierran sin anomalía en ese dominio.

Para exhibir su origen variacional, foliamos
\(M\simeq\mathbb R\times\Sigma\) y escribimos la tríada densitizada y la
conexión de Ashtekar--Barbero
\[
 E_i^a=\frac12\epsilon^{abc}\epsilon_{ijk}e_b^je_c^k,
 \qquad
 A_a^i=\Gamma_a^i+\gamma_{\rm HMT}K_a^i.
\]
La transformada de Legendre de la acción
Palatini--Cartan--Holst especializada produce
\begin{equation}
 H_{\rm tot}
 =\int_\Sigma
 \bigl(\lambda^i\mathcal G_i+N^a\mathcal V_a+N\mathcal H\bigr),
 \label{eq:hamiltoniano-total-gravedad-hmt-rev11}
\end{equation}
donde \(\mathcal G_i\), \(\mathcal V_a\) y \(\mathcal H\) son,
respectivamente, las restricciones de Gauss, difeomorfismos y Hamiltoniana.
La ecuación de Cartan elimina algebraicamente la torsión no propagante; la
curvatura retiene los modos colectivos. Por construcción,
\eqref{eq:hamiltoniano-total-gravedad-hmt-rev11} y las ecuaciones
Euler--Lagrange de la acción covariante poseen el mismo contenido en la
superficie de restricciones.

La relación completa se organiza, por tanto, como
\[
 \text{acción covariante}
 \xleftrightarrow{\ \text{Legendre}\ }
 \text{Hamiltoniano restringido}
 \xrightarrow[\mathscr D]{\ \mathcal Q\ }
 \text{evolución Einstein--Schrödinger}.
\]
La primera flecha es una equivalencia variacional exacta sobre la superficie
clásica de restricciones. La segunda es la realización cuántica definida por
\((\mathcal H_{\rm sec},\mathscr D,\mathcal Q)\); conserva esa equivalencia
cuando se satisfacen la invariancia del dominio y el cierre sin anomalía
declarados arriba.

\begin{exactbox}[title={Resultado variacional y canónico}]
La acción Palatini--Cartan--Holst y el Hamiltoniano total
\eqref{eq:hamiltoniano-total-gravedad-hmt-rev11} describen la misma dinámica
clásica sobre la superficie de restricciones. La cuantización
\(\mathcal Q\) prolonga esa dinámica a
\eqref{eq:einstein-schrodinger-hmt-rev11} sobre el dominio
\(\mathscr D\), con las condiciones explícitas de invariancia y cierre sin
anomalía. Este resultado gobierna la formulación gravitatoria; la
transducción dimensional y la carta decimal de \(G\) se aplican después.
\end{exactbox}

La lectura machiana fija el marco local por restricción de la sección global;
la onda gravitatoria constituye la excitación colectiva de esta geometría,
mientras la torsión algebraica absorbe la respuesta material local.
